\documentclass[10pt,a4paper]{amsart}
\usepackage[margin=19mm]{geometry}
\usepackage{amsmath,amssymb,mathtools}
\usepackage[scr=rsfs,cal=euler]{mathalfa}
\usepackage{xfrac}
\usepackage[unicode,hidelinks,bookmarks=false]{hyperref}
\newcommand{\ie}{i.e.\ }
\newcommand{\cf}{cf.\ }
\newcommand{\resp}{respectively\ }
\newcommand{\Subsection}{\subsection}
\providecommand{\nobreakdash}{\nobreak\hbox{-}}
\setlength{\emergencystretch}{2em}
\multlinegap0pt
\usepackage{bm}
\usepackage{bbm}
\usepackage{dsfont}
\usepackage{xspace}
\usepackage{cancel}
\usepackage{mathtools}
\usepackage{amsthm}
\makeatletter
\@ifundefined{thm}{\newtheorem{thm}{Thm}}{}
\@ifundefined{prop}{\newtheorem{prop}[thm]{Proposition}}{}
\makeatother
\newcommand{\bl}{\textrm{bl}}
\newcommand{\res}{\textrm{res}}
\title[Effective results on projective normality of the first and s]{Effective results on projective normality of the first and second secant varieties}
\author{Doyoung Choi, Jinhyung Park}
\date{Source: arXiv:2606.30002v1 (CC BY 4.0)}
\begin{document}
\maketitle

\noindent\emph{Source and licence.} Effective results on projective normality of the first and second secant varieties. Version arXiv:2606.30002v1 (CC BY 4.0), distributed under the Creative Commons Attribution 4.0 International licence, as recorded in the arXiv licence metadata for this exact version and shown on the abstract page https://arxiv.org/abs/2606.30002v1. Full rights notice: licenses/arxiv-2606.30002-CC-BY-4.0.txt.\par\smallskip

\noindent\textit{Editorial scope.} Counted unit: the Prop whose source label is \texttt{prop:mainvanishing2} in the article (source0.tex lines 470--482, source label \texttt{prop:mainvanishing2}), delivered together with the article's own complete original proof of it (source0.tex lines 485--558, one \texttt{proof} environment of 74 lines). No mathematics is written, repaired or re-derived: statement and proof are the article's own text line for line. Required context retained uncounted: prop:cohvanA (lines 320--329); prop:cohvanB (lines 363--383); prop:mainvanishing1 (lines 401--421). Every definition, display and label of those lines is retained unchanged. Omitted from this package and not redistributed: 1--319, 330--362, 384--400, 422--469, 483--484, 559--993. Nothing inside a retained excerpt is omitted or altered: every retained body is delivered exactly as the article's lines are written, so no omission receipt of any kind accompanies this package. Citations: the retained excerpts carry no citation command at all, so no bibliography entry of the article is transcribed and no reference is redistributed. Reference adaptation: none was needed and none was made; every reference inside the delivered unit resolves inside it, so no unresolved reference or citation is produced. Printed slips kept literal: no printed word, symbol, hypothesis or number of the counted statement or of its proof was altered. Layout and class: the article's own class is \texttt{amsart} with packages including amsmath, amssymb, amsthm, amsfonts, hyperref, cleveref, geometry, blkarray, enumerate, xcolor, graphicx, xy, mathrsfs, eufrak; the wrapper is the lane's pinned \texttt{amsart} editorial wrapper at 10pt on A4, which restates the theorem-like environments the retained text needs and adds the article's own macro definitions that the retained text needs (\texttt{\textbackslash bl}, \texttt{\textbackslash res}), with extra wrapper packages (bm, bbm, dsfont, xspace, cancel, mathtools). The delivered numbering is the unit's own. Assets: no retained span contains a figure environment, a graphics-inclusion command, a PDF-page inclusion, a table or a TikZ picture; no image file is copied into this package, the package declares no asset dependency and no image file is redistributed with it. Licence: version arXiv:2606.30002v1 of this article is distributed under the Creative Commons Attribution 4.0 International licence, as recorded in the arXiv licence metadata for this exact version and shown on the abstract page https://arxiv.org/abs/2606.30002v1. A full rights notice is delivered at licenses/arxiv-2606.30002-CC-BY-4.0.txt. Characters: every retained excerpt is pure ASCII, so the delivered engine, XeLaTeX with its default Latin Modern text font, reports no missing character.
\begin{prop}\label{prop:cohvanA}
Let $k \geq 1$ and $1 \leq i \leq k$ be integers, and $L_i:=\omega_X \otimes H^{m_i} \otimes B_i$ be a line bundle on $X$ with $H$ very ample and $B_i$ nef for an integer $m_i \geq 0$. If $a \geq 1$, $m_1 \geq (k-1)(a+n-1)+1$, and $m_i \geq a+n$ for $2 \leq i \leq k$, then  
$$
H^i(X^k, \mathscr{I}_{\Delta_{1,2} \cup \cdots \cup \Delta_{1,k}}^a \otimes (L_1 \boxtimes \cdots \boxtimes L_k))=0~~\text{ for $i > 0$},
$$
where $\Delta_{1,i}:=\{ (x_1, \ldots, x_k) \in X^k \mid x_1 = x_i\}$ is a pairwise diagonal for $2 \leq i \leq k$. In particular, if $m_1 \geq (k-1)n+1$ and $m_i \geq n+1$ for $2 \leq i \leq k$, then
$$
H^i(X, M_{L_2} \otimes \cdots \otimes M_{L_k} \otimes L_1) = 0~~\text{ for $i>0$}. 
$$
\end{prop}

\begin{prop}\label{prop:cohvanB}
Let $L:=\omega_X \otimes H^m \otimes B$ be a line bundle on $X$ with $H$ very ample and $B$ nef for an integer $m \geq 0$. Then we have the following:
\begin{enumerate}
    \item Let $N$ be a nef line bundle on $X^{[1,2]}$. If $m \geq n+1$, then 
    $$
    H^i(X^{[1,2]}, \rho_{1,2}^* N_{2,L} \otimes N)=0~~\text{ for $i>0$}.
    $$
    In particular, if $m \geq n+2$, then 
    $$
     H^i(X^{[1,2]}, \rho_{1,2}^* A_{2,L} \otimes N)=0~~\text{ for $i>0$}.
    $$
    \item Let $N'$ be a nef line bundle on $X^{[1,2,3]}$, and $N$ be a nef line bundle on $X^{[2,3]}$. If $m  \geq 2n+1$, then 
    $$
    H^i(X^{[1,2,3]}, \rho_{1,2,3}^* N_{3,L} \otimes N')=0~\text{ and }~H^i(X^{[2,3]}, \rho_{2,3}^* N_{3,L} \otimes N)~~\text{ for $i>0$}.
    $$
    In particular, if $m \geq 2n+3$, then 
    $$
    H^i(X^{[1,2,3]}, \rho_{1,2,3}^* A_{3,L} \otimes N')=0~\text{ and }~H^i(X^{[2,3]}, \rho_{2,3}^* A_{3,L} \otimes N)~~\text{ for $i>0$}.
    $$
\end{enumerate}
\end{prop}

\begin{prop}\label{prop:mainvanishing1}
Let $L_1:=\omega_X \otimes H^{m_1} \otimes B_1$ and $L_2:=\omega_X \otimes H^{m_2} \otimes B_2$ be line bundles on $X$ with $H$ very ample and $B_1, B_2$ nef for integers $m_1, m_2 \geq 0$. Then we have the following:
\begin{enumerate}
    \item Let $N$ be a nef line bundle $N$ on $X^{[1,2]}$. If $m_1, m_2 \geq 2n+1$, then 
    $$
    H^i(X^{[1,2]}, \rho_{1,2}^* (M_{2,L_1} \otimes N_{2,L_2}) \otimes N)=0~~\text{ for $i>0$}.
    $$
    In particular, if $m_1 \geq 2n+1$ and $m_2 \geq 2n+2$, then 
     $$
    H^i(X^{[1,2]}, \rho_{1,2}^* (M_{2,L_1} \otimes A_{2,L_2}) \otimes N)=0~~\text{ for $i>0$}.
    $$
    \item Let $N'$ be a nef line bundle on $X^{[1,2,3]}$, and $N$ be a nef line bundle on $X^{[2,3]}$. If $m_1, m_2 \geq 3n+1$, then 
    $$
    H^i(X^{[1,2,3]}, \rho_{1,2,3}^* (M_{3, L_1} \otimes N_{3,L_2}) \otimes N') = 0~\text{and}~H^i(X^{[2,3]}, \rho_{2,3}^* (M_{3, L_1} \otimes N_{3,L_2}) \otimes N)=0~\text{for $i>0$}. 
    $$
    In particular, if $m_1 \geq 3n+1$ and $m_2 \geq 3n+3$, then 
    $$
    H^i(X^{[1,2,3]}, \rho_{1,2,3}^* (M_{3, L_1} \otimes A_{3,L_2}) \otimes N') = 0~\text{and}~H^i(X^{[2,3]}, \rho_{2,3}^* (M_{3, L_1} \otimes A_{3,L_2}) \otimes N)=0~\text{for $i>0$}. 
    $$
\end{enumerate}
\end{prop}

\begin{prop}\label{prop:mainvanishing2}
Let $L_1:=\omega_X \otimes H^{m_1} \otimes B_1$, $L_2:=\omega_X \otimes H^{m_2} \otimes B_2$, and $L_3:=\omega_X \otimes H^{m_3} \otimes B_3$ be line bundles on $X$ with $H$ very ample, $B_1, B_2, B_3$ nef for integers $m_1, m_2, m_3 \geq 0$. Then we have the following:
\begin{enumerate}
\item For a nef line bundle $N$ on $X^{[1,2]}$, if $m_1, m_2, m_3 \geq 3n+2$, then  
$$
H^i(X^{[1,2]}, \rho_{1,2}^* (M_{2,L_1} \otimes M_{2, L_2} \otimes A_{2, L_3}) \otimes N) = 0~~\text{ for $i>0$}. 
$$
\item If $m_1, m_2, m_3 \geq 3n+1$, then 
$$
H^i(X^{[2]}, M_{2,L_1} \otimes M_{2, L_2} \otimes N_{2, L_3}) = 0~~\text{ for $i>0$}. 
$$
\end{enumerate}
\end{prop}

\begin{proof}
$(1)$ Recall from the proof of Proposition \ref{prop:mainvanishing1} $(1)$ that
$$
H^i(X^{[1,2]}, \rho_{1,2}^* (M_{2,L_1} \otimes M_{2, L_2} \otimes A_{2, L_3}) \otimes N)
= H^i(X^{[1,2,3]}, \bl_{1,2,3}^*(L_1 \boxtimes (\rho_{1,2}^* (M_{2, L_2} \otimes A_{2,L_3}) \otimes N))(-F_{1,2}))
$$
for $i>0$. 
We have a short exact sequence
$$
0 \longrightarrow \rho_{1,2,3}^* M_{3, L_2} \longrightarrow \bl_{1,2,3}^* (\mathcal{O}_X \boxtimes \rho_{1,2}^* M_{2, L_2}) \longrightarrow \bl_{1,2,3}^* (L_2 \boxtimes \mathcal{O}_{X^{[1,2]}})(-F_{1,2}) \longrightarrow 0.
$$
It is enough to show that
\begin{align*}
&H^i(X^{[1,2,3]}, \rho_{1,2,3}^* M_{3, L_2} \otimes \bl_{1,2,3}^* (L_1 \boxtimes (\rho_{1,2}^* A_{2,L_3} \otimes N))(-F_{1,2}))=0\\
& H^i(X^{[1,2,3]}, \bl_{1,2,3}^* ( (L_1 \otimes L_2) \boxtimes (\rho_{1,2}^* A_{2,L_3} \otimes N))(-2F_{1,2}))=0.
\end{align*}
Since we may write
\begin{align*}
& \bl_{1,2,3}^* (L_1 \boxtimes (\rho_{1,2}^* A_{2,L_3} \otimes N))(-F_{1,2})\\
&= \rho_{1,2,3}^* N_{3, \omega_X \otimes H^{3n+1}} \otimes \underbrace{\bl_{1,2,3}^* ((H^{m_1 - 3n-1} \otimes B_1) \boxtimes (\rho_{1,2}^* N_{2, H^{m_3-3n-1} \otimes B_3} \otimes N))}_{\text{nef}}\\
& \bl_{1,2,3}^* ( (L_1 \otimes L_2) \boxtimes (\rho_{1,2}^* A_{2,L_3} \otimes N))(-2F_{1,2})\\
& = \rho_{1,2,3}^* A_{3, \omega_X \otimes H^{2n+3}} \otimes \underbrace{\bl_{1,2,3}^* ((H^{m_1 - 2n-3} \otimes B_1 \otimes L_2) \boxtimes (\rho_{1,2}^* T_{2, H^{m_3 - 2n-3} \otimes B_3} \otimes N))}_{\text{nef}},
\end{align*}
the assertion follows from Proposition \ref{prop:mainvanishing1} $(2)$ and Proposition \ref{prop:cohvanB} $(2)$. 


\medskip

\noindent $(2)$ As $L_1$ is very ample and $H^i(X, L_1) = 0$ for $i>0$, we find
$$
R^i \tau_{2,3,*} (\res_{2,3}^* L_1 (-F_2)) = \begin{cases} M_{2, L_1} & \text{for $i=0$} \\ 0 & \text{for $i>0$}. \end{cases}
$$
Then we have
$$
H^i(X^{[2]}, M_{2,L_1} \otimes M_{2, L_2} \otimes N_{2, L_3})
= H^i(X^{[2,3]}, (\res_{2,3}^* L_1 \otimes \tau_{2,3}^* (M_{2, L_2} \otimes N_{2, L_3}))(-F_2))~~\text{ for $i>0$.}
$$
Consider the short exact sequence
$$
0 \longrightarrow \rho_{2,3}^* M_{3, L_2} \longrightarrow \tau_{2,3}^* M_{2, L_2} \longrightarrow \res_{2,3}^* L_2 (-F_2) \longrightarrow 0. 
$$
The desired cohomology vanishing can be deduced from
$$
H^i(X^{[2,3]}, \rho_{2,3}^* M_{3, L_2} \otimes (\res_{2,3}^* L_1 \otimes \tau_{2,3}^* N_{2, L_3})(-F_2))=0~\text{ and }~
H^i(X^{[2,3]}, (\res_{2,3}^* (L_1 \otimes L_2) \otimes \tau_{2,3}^* N_{2, L_3})(-2F_2))=0
$$
for $i>0$. As we may write
$$
(\res_{2,3}^* L_1 \otimes \tau_{2,3}^* N_{2, L_3})(-F_2) = \rho_{2,3}^* N_{3, \omega_X \otimes H^{3n+1}} \otimes \underbrace{\res_{2,3}^* (H^{m_1 - 3n-1} \otimes B_1) \otimes \tau_{2,3}^* T_{2, H^{m_3-3n-1} \otimes B_3}}_{\text{nef}},
$$
the first holds by Proposition \ref{prop:mainvanishing1} $(2)$. For the second, notice that $\rho_{1,2,3,*}' \mathcal{O}_{X^{[1,2,3]}} = \mathcal{O}_{X^{[2,3]}} \oplus \tau_{2,3}^* \mathcal{O}_{X^{[2]}}(-\delta_2)$. Thus it suffices to show that
$$
H^i(X^{[1,2,3]}, \res_{1,2,3}^* (L_1 \otimes L_2) \otimes \tau_{1,2,3}^* T_{2, L_3}(-2F_{1,2}))=0~~\text{ for $i>0$}.
$$
Recall that $X^{[1,2,3]}$ is obtained from $X^3$ by the composition of blow-ups
$$
b \colon X^{[1,2,3]} \xrightarrow{~\bl_{1,2,3}~} X \times X^{[1,2]} \xrightarrow{~\operatorname{id}_X \times \bl_{1,2}~} X^3.
$$
We may write
$$
\res_{1,2,3}^* (L_1 \otimes L_2) \otimes \tau_{1,2,3}^* T_{2, L_3}(-2F_{1,2})
= b^* ((L_1 \otimes L_2) \boxtimes L_3 \boxtimes L_3)(-2F_{1,2}).
$$
Note that 
$$
b_* \mathcal{O}_{X^{[1,2,3]}}(-2F_{1,2}) = (\operatorname{id}_X \times \bl_{1,2})_* \mathscr{I}_{\mathcal{W}_{1,2}}^2 = \mathscr{I}_{\Delta_{1,2} \cup \Delta_{1,3}}^2. 
$$
Then 
$$
H^i(X^{[1,2,3]}, \res_{1,2,3}^* (L_1 \otimes L_2) \otimes \tau_{1,2,3}^* T_{2, L_3}(-2F_{1,2}))
= H^i(X^3, \mathscr{I}_{\Delta_{1,2} \cup \Delta_{1,3}}^2 \otimes ((L_1 \otimes L_2) \boxtimes L_3 \boxtimes L_3))=0~~\text{ for $i>0$}
$$
by Proposition \ref{prop:cohvanA}.
\end{proof}

\end{document}
